% s-301-p301-m002; original PDF 301, 302
\Aphorism{s-301-p301-m002}{357}
Mein Dasein \Foreign{a posteriori} rechtfertigen — das ist Kleinkram; aber mein Dasein als intelligiblen freien Willensact, als Schuld ansetzen dürfen, Schuld nicht nur an mir selbst, sondern „an Allem und für Alle“, um mit einem Neurastheniker aus Dostojewskijs Klinik zu reden, Schuld am Dasein der Welt, Schuld am qualvollen Sterben des Laurentius oder irgendwelcher Mücke auf irgendwelchem Planeten des Sirius — das ist Moral grossen Stiles, leider zu grossen für uns tellurische Däumlinge. Schopenhauers Ethik ist eigentlich ein genialer Versuch, durch den Universalismus zum Nihilismus in der Moral zu kommen. Man erwäge einen Satz wie den folgenden (Parerga II, 165): „Wer etwas tiefer zu denken fähig ist wird bald absehn, dass die menschlichen Begierden nicht erst auf dem Punkte anfangen können, sündlich zu sein, wo sie, in ihren individuellen Richtungen einander zufällig durchkreuzend, Übel von der einen und Böses von der andern Seite veranlassen; sondern dass, wenn dieses ist, sie auch schon ursprünglich und ihrem Wesen nach sündlich und verwerflich sein müssen, folglich der ganze Wille zum Leben selbst ein verwerflicher ist.“ Das ist der verfängliche Umweg über das All zum Nichts: die einzelne Sünde geht unter in der allgemeinen Sündhaftigkeit — ein drolliger Streich wie im Märchen, wo der Übelthäter seine Verfolger irreführt, indem er alle Hausthüren des Quartiers mit Kreide anstreicht. Schopenhauer, der absolute Moralist und erste Immoralist, hat uns den Geschmack für jede specielle Moral mit begrenztem Pflichtenkreis verdorben, so wie ihn früher das Christenthum verdarb, sowie ihn jede amphibische Weltanschauung verdirbt, die in den Schranken des practischen Handelns die erkannte Unendlichkeit und Transcendenz des Weltganzen zum Ausdruck bringen will. Die Moral verträgt keine Verankerung jenseits der Welt, nicht einmal jenseits der Erde oder des Menschen; insofern aber die grössten Moralphilosophen diesen Irrweg, der zum moralischen Indifferentismus führt, betreten haben, darf man mit gutem Recht sagen: wenn die Moral heute noch steht, dann \Sperr{trotz }ihren Begründern\Sperr{!}
\par

